
\documentclass[11pt, letterpaper]{nyjm}

\usepackage{amsmath, amssymb, amsthm}
\usepackage{hyperref}
\usepackage{geometry}
\geometry{margin=1in}

% Theorem environments
\newtheorem{theorem}{Theorem}[section]
\newtheorem{lemma}[theorem]{Lemma}
\newtheorem{corollary}[theorem]{Corollary}
\newtheorem{remark}[theorem]{Remark}
\newtheorem{definition}{Definition}
\newtheorem{proposition}{Proposition}

% Custom commands
\newcommand{\R}{\mathbb{R}}
\newcommand{\Sn}{S^{n-1}}
\newcommand{\Tdo}{T^{\delta}_{\omega}}
\newcommand{\Sph}{\mathbb{S}^{n-1}}
\newcommand{\calP}{\mathcal{P}}
\newcommand{\range}{\operatorname{range}}
\newcommand{\rep}{\operatorname{rep}}
\newcommand{\LowerIncidence}{\operatorname{LowerIncidence}}

\title{The Kakeya Maximal Function Conjecture and the Dyadic Bound}

\author{Johan Aspegren}
\email{jaspegren@outlook.com}
\date{October 29, 2026}

\subjclass[2020]{Primary 42B37; Secondary 28A75}
\keywords{Kakeya Conjecture, Harmonic Analysis, Incidence Geometry}

\begin{document}
	
	\begin{abstract}
		We use a dyadic $L_2$-bound to prove the Kakeya maximal function conjecture. We study the Kakeya maximal function conjecture through a dyadic \(L^2\) reduction and localized incidence counting. Combining greedy tube selection, Hall's matching theorem, and geometric packing estimates, we reduce the conjectured bound to a localized mass comparison and an incidence counting estimate.  
	\end{abstract}
	
	\maketitle
	
	\tableofcontents
	
	\section{Introduction}
	%------------------------------------------------------------
	
	We define a line $l_i$ as
	\[
	l_i := \{ y \in \R^n \mid \exists t \in \R,\ y = a_i + t\omega_i \},
	\]
	where $\omega_i \in S^{n-1}.$
	We define the $\delta$-tubes as $\delta$-neighborhoods of lines on $B(0,1)$ that
	intersect the ball $B(0, 1/2)$:
	\[
	T^{\delta}_i := \{ x \in B(0,1) \mid |x - y| < \delta,\ y \in l_i
	\text{ and } B(0,1/2) \cap l_i \neq \emptyset \}.
	\]
	So that lengths of our $\delta$-tubes are bounded also from below. The
	\emph{order of intersection} is defined as the number of tubes intersecting in an
	intersection. We define $A \lesssim B$ to mean that there exists a constant $C_n$
	depending only on $n$ such that $A \leq C_n B$. We say that tubes are
	$\delta$-\emph{separated} if their angles are $\delta$-separated. Moreover, let
	$f \in L^1_{\mathrm{loc}}(\R^n)$. For each tube in $B(0,1)$ define $a$ as its
	center of mass. Define the \emph{Kakeya maximal function}
	\[
	f^*_\delta : S^{n-1} \to \R \quad \text{via} \quad
	f^*_\delta(\omega) = \sup_{a \in \R^n} \frac{1}{|T^{\delta}_\omega(a)|}
	\int_{T^{\delta}_\omega(a)} |f(y)|\, dy.
	\]
	In this paper any constant can depend on dimension $n$. In the study of the Kakeya
	maximal function conjecture we are aiming at the following bounds:
	\begin{equation}\label{eq:maxfunc}
		\| f^*_\delta \|_p \leq C_\epsilon \delta^{-n/p + 1 - \epsilon} \| f \|_p,
	\end{equation}
	for all $\epsilon > 0$ and some $n \leq p \leq \infty$. A very important
	reformulation of the problem by Tao is the following. A bound of the form
	\eqref{eq:maxfunc} follows from a bound of the form
	\begin{equation}\label{eq:tao}
		\left\| \sum_{\omega \in \Omega} \mathbf{1}_{T_\omega(a_\omega)}(x) \right\|_{p/(p-1)}
		\leq C_\epsilon \delta^{-n/p + 1 - \epsilon},
	\end{equation}
	for all $\epsilon > 0$, and for any set of $N \leq \delta^{1-n}$ $\delta$-separated
	$\delta$-tubes. See for example \cite{Kroc} or \cite{Tao99}. It is enough to
	consider the case $p = n$ and the rest of the cases will follow via interpolation
	\cite{Kroc, Tao99}. In this paper any constant can depend on dimension $n$. Our
	main theorem is the following.
	
	\begin{theorem}\label{thm:main}
		Let there be $N \sim \delta^{1-n}$ $\delta$-tubes that are $\delta$-separated.
		Then we have that
		\[
		\left\| \sum_{\omega \in \Omega} \mathbf{1}_{T_\omega(a_\omega)}(x)
		\right\|_{n/(n-1)} \lesssim (\ln N)^{(n-1)/n}.
		\]
	\end{theorem}
	
	The case $n = 2$ of the Kakeya maximal function conjecture is well known to be
	true \cite{Kroc}. The case $n = 1$ is trivial.
	A formalization project for this paper in Lean 4 is available at
	\cite{Aspegren2026Lean}.
	%------------------------------------------------------------
	\section{Previously Known Results}
	%------------------------------------------------------------
	
	We will use the following bound for the pairwise intersections of $\delta$-tubes.
	
	\begin{lemma}[C\'ordoba]\label{lem:cordoba}
		For any pair of directions $\omega_i, \omega_j \in S^{n-1}$ and any pair of
		points $a, b \in \R^n \cap B(0,1)$, we have
		\[
		|T^{\delta}_{\omega_i}(a) \cap T^{\delta}_{\omega_j}(b)|
		\lesssim \frac{\delta^n}{|\omega_i - \omega_j|}.
		\]
	\end{lemma}
	
	A proof can be found for example in \cite{Kroc}. For any (spherical) cap
	$\Omega \subset S^{n-1}$ with $|\Omega| \gtrsim \delta^{n-1}$, $\delta > 0$,
	define its $\delta$-\emph{entropy} $N_\delta(\Omega)$ as the maximum possible
	cardinality of a $\delta$-separated subset of $\Omega$.
	
	\begin{lemma}\label{lem:entropy}
		In the notation just defined,
		\[
		N_\delta(\Omega) \sim \frac{|\Omega|}{\delta^{n-1}}.
		\]
	\end{lemma}
	
	Again, a proof can essentially be found in \cite{Kroc}.
	
	We also need the following lemma.
	\begin{lemma}[Vitali covering lemma for balls of equal radius]
		\label{lem:vitali-equal}
		Let $\rho>0$ and let $\{B(x_\alpha,\rho)\}_{\alpha\in\mathcal A}$ be a family of balls in $\mathbb R^n$
		of the same radius (all open or all closed). Then there is $\mathcal S\subseteq\mathcal A$ such that
		\begin{enumerate}
			\item[(i)] the balls $B(x_\alpha,\rho)$, $\alpha\in\mathcal S$, are pairwise disjoint,
			\item[(ii)] $\mathcal S$ is countable, and
			\item[(iii)] for every $\alpha\in\mathcal A$ there is $\beta\in\mathcal S$ with
			$|x_\alpha-x_\beta|\le2\rho$ (in particular $B(x_\alpha,\rho)\subseteq B(x_\beta,3\rho)$). Hence
			\[
			\bigcup_{\alpha\in\mathcal A}B(x_\alpha,\rho)\subseteq\bigcup_{\beta\in\mathcal S}B(x_\beta,3\rho),
			\qquad
			\Bigl|\bigcup_{\alpha\in\mathcal A}B(x_\alpha,\rho)\Bigr|\le3^n\sum_{\beta\in\mathcal S}|B(x_\beta,\rho)|.
			\]
		\end{enumerate}
		If $\mathcal A$ is finite, $\mathcal S$ can be taken to be any maximal pairwise disjoint subfamily.
	\end{lemma}
	
	\begin{proof}
		By Zorn's lemma there is a maximal pairwise disjoint subfamily $\mathcal S\subseteq\mathcal A$
		(the union of a chain of pairwise disjoint subfamilies is pairwise disjoint, so chains have upper
		bounds). For finite $\mathcal A$ this is just a greedy choice. This gives (i).
		
		(ii) Each ball has nonempty interior, so it contains a point of $\mathbb Q^n$ in its interior.
		Disjoint balls contain distinct such points, so $\mathcal S$ injects into $\mathbb Q^n$.
		
		(iii) Let $\alpha\in\mathcal A$. If $\alpha\in\mathcal S$ take $\beta=\alpha$. Otherwise, by maximality of
		$\mathcal S$, $B(x_\alpha,\rho)$ meets some $B(x_\beta,\rho)$ with $\beta\in\mathcal S$. Pick $y$ in the
		intersection. Then
		\[
		|x_\alpha-x_\beta|\le|x_\alpha-y|+|y-x_\beta|\le2\rho,
		\]
		and for $z\in B(x_\alpha,\rho)$,
		\[
		|z-x_\beta|\le|z-x_\alpha|+|x_\alpha-x_\beta|\le\rho+2\rho=3\rho .
		\]
		For open balls the inequalities are strict in the first step and the conclusion holds with the
		open ball $B(x_\beta,3\rho)$; for closed balls with the closed one. The measure bound follows from
		countable subadditivity and $|B(x_\beta,3\rho)|=3^n|B(x_\beta,\rho)|$.
	\end{proof}
	%------------------------------------------------------------
	\section{The Generalized Lemma}
	%-------------------------------------------------------
	
	For the Kakeya set define the  dyadic level sets as 
	\[
	F_m := \left\{ x \in \R^n \;\middle|\; 2^{m} \leq \sum_{i=1}^{N}
	\mathbf{1}_{T_i}(x) < 2^{m+1} \right\}.
	\]
	
	\begin{lemma}[Generalized C\'ordoba lemma] \label{lem:corgeneralized}
		We have 
		
		\[
		\bigcap_{i=1}^{\sim 2^{k}} T_i^{\delta} \subset B(a,C2^{-k/(n-1)})  
		\]
	\end{lemma}
	
	\begin{proof}
		Let us suppose that
		\[
		2^k=\delta^{-\beta},
		\qquad 0<\beta\leq n-1,
		\]
		and let us suppose that a tube $T_j$ intersecting
		$T_i\cap E_{2^k}$ has its direction outside of a cap of size
		$\sim\delta^{n-1-\beta}$ on the unit sphere. Then the angle between
		$T_j$ and $T_i$ is greater than $\sim\delta^{1-\beta/(n-1)}$.
		Thus, the intersection 
		$\bigcap_{i=1}^{2^k}T_i \subset T_i \cap T_j \subset B(a,C2^{-k/(n-1)}).$ 
		So we can suppose that the directions in the intersection
		$E_{2^k}\cap T_i\cap T_j$ belong to a cap of size
		$\sim\delta^{n-1-\beta}$. However, the cap contains $2^k$ tube
		directions. 
		Thus, for any tube $T_i$ in the intersection, there exists
		a tube $T_j$ such that the angle between $T_i$ and $T_j$ is
		$\sim\delta^{1-\beta/(n-1)}$ and again \\ $\bigcap_{i=1}^{2^k}T_i \subset T_i \cap T_j \subset B(a,C2^{-k/(n-1)})$ is valid.
		Thus we proved the lemma \ref{lem:corgeneralized}.
		
		
	\end{proof}
	
	
	
	\section{The Proof}
	
	
	Now, 
	\begin{equation} \label{eq:final1}
		\begin{aligned}	
			&\left\| \sum_{i=1}^{N} \mathbf{1}_{T_i(a_i)}(x)
			\right\|_{n/(n-1)}^{n/(n-1)} \\
			&\sim \int_{B(0,1)}(\sum_{i=1}^{N} \mathbf{1}_{T_i(a_i)}(x))^{n/(n-1)} \\
			&\sim \sum_m \int_{F_m}(\sum_{i=1}^{N} \mathbf{1}_{T_i(a_i)}(x))^{n/(n-1)} \\
			&\sim \sum_m 2^{-m(2-n/(n-1))}\int_{F_m}(\sum_{i=1}^{N} \mathbf{1}_{T_i(a_i)}(x))^{2} \\
			&\sim \sum_m 2^{-m(2-n/(n-1))}\sum_{j=1}^{N}\sum_{i=1}^{N}|T^\delta_i(a_i) \cap T^\delta_j(a_j) \cap F_m|, 
		\end{aligned}
	\end{equation} 
	where we used the dyadic decomposition.
	Moreover, when looking at the last term inside the parentheses in \eqref{eq:final1} we have the dyadic $L_2$-bound 
	\begin{equation} \label{eq:final2}
		\begin{aligned}	
			&\sum_{j=1}^{N}\sum_{i=1}^{N}|T^\delta_i(a_i) \cap T^\delta_j(a_j) \cap F_m| \\
			& = \sum_{j=1}^{N}\sum_{i=1}^{N} \sum_{l}|T^\delta_i(a_i) \cap T^\delta_j(a_j)\cap B(b_{l},C2^{-m/(n-1)}) \cap F_{ml}|
			,
		\end{aligned}
	\end{equation}
	where we used the generalized lemma of C{\'o}rdoba \ref{lem:corgeneralized}, and that the intersection is unique if it exists, and included in a ball $B(b_{ij},C2^{-m/(n-1)}).$ Moreover we used that   
	\[
	F_m = \bigcup_l\bigcap_{k=1}^{\sim 2^m} T_{k,l}^{\delta}, 
	\]
	where the partition is disjoint.
	Each intersection center of  $\bigcap_{k=1}^{\sim 2^m} T_{k,l}^{\delta}$ was chosen as the center point $b_{l},$ and each $F_{ml}$ represents an individual intersection $\bigcap_{k=1}^{\sim 2^m} T_{k,l}^{\delta}.$ 
	
	\begin{theorem}[Greedy tube selection and averaging]
		\label{thm:greedy-selection}
		
		Let $q=2^m$, and let
		$T_1^\delta,\ldots,T_N^\delta$ be the original
		$\delta$-tubes. Define the dyadic multiplicity
		level set
		\[
		F_m=
		\left\{
		x\in\mathbb R^n:
		q\leq\sum_{i=1}^N
		\mathbf{1}_{T_i^\delta}(x)<2q
		\right\}.
		\]
		
		Assume that $|F_m|>0$.
		
		Then there exist distinct tube indices
		$j_1,\ldots,j_K$, where
		\[
		K=\left\lceil\frac Nq\right\rceil,
		\qquad Kq\leq 2N,
		\]
		such that
		\[
		\boxed{
			|F_m|\lesssim
			\sum_{k=1}^K
			|T_{j_k}^\delta\cap F_m|.
		}
		\]
		
		Moreover, the tubes can be selected greedily
		so that their newly covered pieces
		\[
		Q_k:=
		T_{j_k}^\delta\cap
		\left(
		F_m\setminus
		\bigcup_{r=1}^{k-1}T_{j_r}^\delta
		\right)
		\]
		are pairwise disjoint and satisfy
		\[
		|Q_1|\geq |Q_2|\geq\cdots\geq |Q_K|.
		\]
		
		In fact, the stronger estimate
		\[
		\boxed{
			|F_m|\leq
			\frac{1}{1-e^{-1}}
			\sum_{k=1}^K|Q_k|
		}
		\]
		holds.
		
	\end{theorem}
	
	\begin{proof}
		
		Set
		\[
		F_m^{(0)}=F_m.
		\]
		
		For each $k\geq1$, choose $j_k$ so that
		\[
		|T_{j_k}^\delta\cap F_m^{(k-1)}|
		=
		\max_{1\leq i\leq N}
		|T_i^\delta\cap F_m^{(k-1)}|.
		\]
		
		Define
		\[
		Q_k=T_{j_k}^\delta\cap F_m^{(k-1)}
		\]
		and
		\[
		F_m^{(k)}
		=
		F_m^{(k-1)}\setminus Q_k.
		\]
		
		The sets $Q_k$ are pairwise disjoint
		by construction.
		
		Since
		\[
		F_m^{(k)}\subseteq F_m^{(k-1)},
		\]
		we have
		\[
		\begin{aligned}
			|Q_{k+1}|
			&=
			\max_{1\leq i\leq N}
			|T_i^\delta\cap F_m^{(k)}|\\
			&\leq
			\max_{1\leq i\leq N}
			|T_i^\delta\cap F_m^{(k-1)}|\\
			&=|Q_k|.
		\end{aligned}
		\]
		
		Thus,
		\[
		|Q_1|\geq|Q_2|\geq\cdots.
		\]
		
		Next, by the dyadic multiplicity condition,
		\[
		q\leq
		\sum_{i=1}^N\mathbf{1}_{T_i^\delta}(x)
		<2q
		\qquad (x\in F_m).
		\]
		
		Consequently,
		\[
		\begin{aligned}
			N|Q_k|
			&=
			N\max_{1\leq i\leq N}
			|T_i^\delta\cap F_m^{(k-1)}|\\
			&\geq
			\sum_{i=1}^N
			|T_i^\delta\cap F_m^{(k-1)}|\\
			&=
			\int_{F_m^{(k-1)}}
			\sum_{i=1}^N
			\mathbf{1}_{T_i^\delta}(x)\,dx\\
			&\geq q|F_m^{(k-1)}|.
		\end{aligned}
		\]
		
		Hence,
		\[
		\boxed{
			|Q_k|\geq
			\frac qN|F_m^{(k-1)}|.
		}
		\]
		
		It follows that
		\[
		\begin{aligned}
			|F_m^{(k)}|
			&=
			|F_m^{(k-1)}|-|Q_k|\\
			&\leq
			\left(1-\frac qN\right)
			|F_m^{(k-1)}|.
		\end{aligned}
		\]
		
		Iterating gives
		\[
		\begin{aligned}
			|F_m^{(k)}|
			&\leq
			\left(1-\frac qN\right)^k|F_m|\\
			&\leq e^{-kq/N}|F_m|.
		\end{aligned}
		\]
		
		Choose
		\[
		K=\left\lceil\frac Nq\right\rceil.
		\]
		
		Since $F_m$ is nonempty and $q\leq N$,
		\[
		Kq\leq N+q\leq 2N.
		\]
		
		Furthermore,
		\[
		|F_m^{(K)}|
		\leq e^{-Kq/N}|F_m|
		\leq e^{-1}|F_m|.
		\]
		
		Since the pieces $Q_k$ are disjoint,
		\[
		\begin{aligned}
			\sum_{k=1}^K|Q_k|
			&=
			|F_m|-|F_m^{(K)}|\\
			&\geq
			(1-e^{-1})|F_m|.
		\end{aligned}
		\]
		
		Therefore,
		\[
		\boxed{
			|F_m|
			\leq
			\frac{1}{1-e^{-1}}
			\sum_{k=1}^K|Q_k|.
		}
		\]
		
		Finally, since
		\[
		Q_k\subseteq T_{j_k}^\delta\cap F_m,
		\]
		we obtain
		\[
		\begin{aligned}
			|F_m|
			&\lesssim\sum_{k=1}^K|Q_k|\\
			&\leq
			\sum_{k=1}^K
			|T_{j_k}^\delta\cap F_m|.
		\end{aligned}
		\]
		
		This proves the theorem.
	\end{proof}
	
	Next,
	\begin{equation} \label{eq:final3}
		\begin{aligned}	
			&(2^m)^2|F_m| \sim \sum_{j=1}^{N}\sum_{i=1}^{N}\sum_{l}|T^\delta_i(a_i) \cap T^\delta_j(a_j)\cap B(b_{l},C2^{-m/(n-1)}) \cap F_{ml}| \\ 
			&\sim 2^m\sum_{i=1}^{N}\sum_{l=1}^{\alpha_i}|T^\delta_i(a_i) \cap B(b_{il},C2^{-m/(n-1)}) \cap F_{ml}| \\
			&\sim (2^m)^2\sum_{i=1}^{K}\sum_{l=1}^{\alpha_i}|T^\delta_i(a_i) \cap B(b_{il},C2^{-m/(n-1)}) \cap F_{ml} \cap G_i| \\
			&\sim  (2^{m})^2\sum_{j=1}^{M}|T^\delta_j(a_j) \cap B(b_{j},3C2^{-m/(n-1)}) \cap F_m \cap G_j|,		
		\end{aligned}
	\end{equation}
	where $K$ pieces are chosen as in \ref{thm:greedy-selection}. However, while  re-indexing in the last sum, \emph{a priori} we may need to choose some tubes multiple times.
	However we can represent each piece in the important cases with distinct $\delta$-tubes. 
	\begin{theorem}[Distinct tube representatives via Hall's theorem]
		\label{thm:hall-representatives}
		
		Let $n\ge 3$, $q=2^m$, and let
		$T_1^\delta,\ldots,T_N^\delta$ be the original tubes.
		
		Let $\mathcal P$ be a finite family of nonempty
		localized pieces contained in the dyadic level set $F_m$.
		
		
		For each $P\in\mathcal P$, define
		\[
		\mathcal N(P)
		:=
		\{r\in\{1,\ldots,N\}:
		P\cap T_r^\delta\neq\emptyset\}.
		\]
		
		Assume that each piece intersects at least $q$ original
		tubes, so that
		\[
		|\mathcal N(P)|\ge q
		\qquad(P\in\mathcal P).
		\]
		
		Suppose further that each original tube intersects
		at most $Cq^{1/(n-1)}$ distinct pieces:
		\[
		\#\{P\in\mathcal P:
		P\cap T_r^\delta\neq\emptyset\}
		\le Cq^{1/(n-1)}
		\qquad(1\le r\le N).
		\]
		
		If
		\[
		q^{(n-2)/(n-1)}\ge C,
		\]
		then there exists an injective map
		\[
		\sigma:\mathcal P\longrightarrow\{1,\ldots,N\}
		\]
		such that
		\[
		P\cap T_{\sigma(P)}^\delta\neq\emptyset
		\qquad(P\in\mathcal P).
		\]
		In particular, distinct localized pieces admit
		distinct original tube representatives.
	\end{theorem}
	
	\begin{proof}
		Let $\mathcal S\subseteq\mathcal P$ be nonempty and define
		\[
		\mathcal N(\mathcal S)
		:=
		\bigcup_{P\in\mathcal S}\mathcal N(P).
		\]
		
		Double counting the incidences between the pieces
		in $\mathcal S$ and the original tubes gives
		\[
		\begin{aligned}
			q|\mathcal S|
			&\le \sum_{P\in\mathcal S}|\mathcal N(P)|\\
			&=
			\sum_{r\in\mathcal N(\mathcal S)}
			\#\{P\in\mathcal S:
			P\cap T_r^\delta\neq\emptyset\}\\
			&\le
			Cq^{1/(n-1)}|\mathcal N(\mathcal S)|.
		\end{aligned}
		\]
		
		Consequently,
		\[
		\boxed{
			|\mathcal N(\mathcal S)|
			\ge
			\frac{q^{(n-2)/(n-1)}}{C}
			|\mathcal S|.
		}
		\]
		
		Since $q^{(n-2)/(n-1)}\ge C$, we obtain
		\[
		|\mathcal N(\mathcal S)|\ge|\mathcal S|.
		\]
		
		The same inequality is trivial for the empty set.
		Thus Hall's marriage condition holds for every
		$\mathcal S\subseteq\mathcal P$.
		
		By Hall's marriage theorem, there exists an
		injective selection
		\[
		\sigma:\mathcal P\longrightarrow\{1,\ldots,N\},
		\qquad
		\sigma(P)\in\mathcal N(P).
		\]
		
		Hence
		\[
		P\cap T_{\sigma(P)}^\delta\neq\emptyset
		\]
		for every $P\in\mathcal P$.
	\end{proof}
	\begin{corollary}[Intersection-compatible bijection fixing selected tubes]
		\label{cor:restricted-intersection-bijection}
		
		Let $B\subseteq\{1,\ldots,N\}$ be the selected tube
		indices, and let $\mathcal P$ be a finite family of
		nonempty localized trimmed pieces.
		
		For each $P\in\mathcal P$, define
		\[
		\mathcal N_B(P)
		:=
		\{r\in\{1,\ldots,N\}\setminus B:
		P\cap T_r^\delta\ne\varnothing\}.
		\]
		
		Suppose that, for some constants $a\ge b>0$,
		\[
		|\mathcal N_B(P)|\ge a
		\qquad(P\in\mathcal P),
		\]
		and
		\[
		\#\{P\in\mathcal P:
		P\cap T_r^\delta\ne\varnothing\}
		\le b
		\qquad(r\notin B).
		\]
		
		Then there exists an injective map
		\[
		\sigma:\mathcal P
		\longrightarrow \{1,\ldots,N\}\setminus B
		\]
		such that
		\[
		P\cap T_{\sigma(P)}^\delta\ne\varnothing
		\qquad(P\in\mathcal P).
		\]
		
		Consequently, setting
		\[
		D:=B\cup\sigma(\mathcal P),
		\]
		there exists a bijection
		\[
		\Phi:B\sqcup\mathcal P\longrightarrow D
		\]
		fixing every index in $B$, with
		\[
		|D|=|B|+|\mathcal P|.
		\]
		
	\end{corollary}
	
	\begin{proof}
		For every nonempty $S\subseteq\mathcal P$,
		double counting the incidences gives
		\[
		a|S|
		\le \sum_{P\in S}|\mathcal N_B(P)|
		\le b\left|
		\bigcup_{P\in S}\mathcal N_B(P)
		\right|.
		\]
		
		Since $a\ge b>0$, it follows that
		\[
		\left|
		\bigcup_{P\in S}\mathcal N_B(P)
		\right|
		\ge |S|.
		\]
		
		Thus Hall's marriage condition holds, and
		there exists an injective map
		\[
		\sigma:\mathcal P
		\longrightarrow\{1,\ldots,N\}\setminus B
		\]
		satisfying
		\[
		\sigma(P)\in\mathcal N_B(P)
		\qquad(P\in\mathcal P).
		\]
		
		In particular,
		\[
		P\cap T_{\sigma(P)}^\delta\ne\varnothing.
		\]
		
		Because $\sigma$ is injective and
		$\sigma(\mathcal P)\cap B=\varnothing$,
		the map
		\[
		\Phi:B\sqcup\mathcal P\longrightarrow D
		\]
		defined by
		\[
		\Phi(\ell)=\ell\quad(\ell\in B),
		\qquad
		\Phi(P)=\sigma(P)\quad(P\in\mathcal P)
		\]
		is bijective.
		
		Consequently,
		\[
		|D|=|B|+|\mathcal P|.
		\]
	\end{proof}
	
	
	The argument for the upper $b$ bound is that you can cover the unit tube with very few balls of radius $C2^{-m/(n-1)}$: 
	\[
	\max_i \alpha_i \leq C^{-1}2^{m/(n-1)} \leq c2^{m},
	\]
	but we need to track constants.
	\begin{lemma}[Disjoint balls meeting a tube] \label{lem:disjoint-ball-tube}
		Let $n\ge2$, let $I\subset\mathbb{R}^n$ be a line segment
		of length $L$, and define
		\[
		T^\delta=\{x\in\mathbb{R}^n:\operatorname{dist}(x,I)<\delta\}.
		\]
		Let $B_1,\ldots,B_M$ be pairwise disjoint balls of radius $R$,
		where $0<\delta\le R$. Then
		\[
		\#\{j:T^\delta\cap B_j\ne\varnothing\}
		\le K_n\left(1+\frac{L}{R}\right),
		\]
		where $K_n$ depends only on the dimension.
		In particular, if $L=1$ and $R\le1$, then
		\[
		\#\{j:T^\delta\cap B_j\ne\varnothing\}\le \frac{2K_n}{R}.
		\]
	\end{lemma}
	
	\begin{proof}
		Let $\mathcal{J}=\{j:T^\delta\cap B_j\ne\varnothing\}$.
		For each $j\in\mathcal{J}$, choose $x_j\in T^\delta\cap B_j$.
		If $y\in B_j$, then
		\[
		\operatorname{dist}(y,I)
		\le |y-x_j|+\operatorname{dist}(x_j,I)
		<2R+\delta\le3R.
		\]
		Consequently,
		\[
		\bigcup_{j\in\mathcal{J}}B_j
		\subseteq\{x:\operatorname{dist}(x,I)<3R\}.
		\]
		The set on the right is contained in a cylinder of length
		$L+6R$ and cross-sectional radius $3R$.
		Writing $\omega_d$ for the volume of the unit ball in
		$\mathbb{R}^d$, disjointness gives
		\[
		\#\mathcal{J}\,\omega_n R^n
		\le \omega_{n-1}(3R)^{n-1}(L+6R).
		\]
		Dividing by $\omega_n R^n$ proves the result.
	\end{proof}
	
	So we have the both conditions for the lemma for the relevant scales.
	We prove the last equation of\eqref{eq:final3}.
	According to greedy-selection \ref{thm:greedy-selection} we have the following cover of a major portion of $F_m$: 
	\[
	\begin{aligned}
		&\bigcup_{\ell\in B}\bigcup_{i_\ell}
		T_\ell^\delta
		\cap B(b_{\ell i_\ell},R)
		\cap F_m
		\cap G_\ell\\
		&\quad\subseteq
		\bigcup_{\ell\in B}\bigcup_{j=1}^{m_\ell}
		T_\ell^\delta
		\cap B(b_{\ell, j},3R)
		\cap F_m
		\cap G_\ell\\
		&\quad=
		\bigcup_{k\in D}
		T_{\ell(k)}^\delta
		\cap B(b_{k(\ell,j)},3R)
		\cap G_{\ell(k)}
		\cap F_m,
	\end{aligned}
	\]
	where
	\[
	R=Cq^{-1/(n-1)},\qquad q=2^m,
	\]
	and $m_\ell$ denotes the number of balls selected by
	the Vitali covering theorem for the tube $\ell\in B$.
	The set $D$ is obtained by bijectively re-indexing the
	selected tube-ball pairs, when the suitable bijection exists.
	For the existence conditions for suitable bijection we need the following lemma.
	We first define 
	\[
	\mathcal A
	:=
	\{(\ell,j):\ell\in B',\ 1\le j\le m_\ell\}
	\]
	be the set of selected tube-ball pairs.
	Define a bijective re-indexing for the Vitali- balls
	\[
	k:\mathcal A\longrightarrow D,
	\qquad
	(\ell,j)\longmapsto k(\ell,j),
	\]
	with inverse
	\[
	k^{-1}(k_0)
	=
	(\ell(k_0),j(k_0)),
	\qquad k_0\in D.
	\]
	In particular,
	\[
	b_{k(\ell,j)}=b_{\ell,j},
	\qquad
	b_{k_0}=b_{\ell(k_0),j(k_0)}.
	\]
	
	
	\begin{lemma}[Packing of exceptional Vitali balls]
		\label{lem:exceptional-vitali}
		
		Let $n\ge 3$, $N\asymp\delta^{1-n}$,
		and let $T_1^\delta,\ldots,T_N^\delta$
		have pairwise $\delta$-separated directions.
		
		Let $q=2^m\le N$, and set
		\[
		R=Aq^{-1/(n-1)},\qquad
		\delta\le R\le 1,
		\]
		where $A>0$ is fixed.
		
		Let
		\[
		\mathcal V=\{V_j=B(b_j,R):j\in J\}
		\]
		be a finite family of pairwise disjoint
		Vitali balls meeting $F_m$.
		
		For each $j$ in a subfamily
		$J_{\mathrm{exc}}\subseteq J$, suppose
		there is a nonempty piece
		$P_j\subseteq V_j\cap F_m$ satisfying
		\[
		d(P_j):=
		\#\{r\le N:P_j\cap T_r^\delta
		\ne\varnothing\}>\frac N2.
		\]
		
		Then
		\[
		\boxed{|J_{\mathrm{exc}}|\le C_{n,A},}
		\]
		where the constant also depends on
		the fixed constants in
		$N\asymp\delta^{1-n}$.
	\end{lemma}
	
	\begin{proof}
		Write
		\[
		I_j=\{r\le N:P_j\cap T_r^\delta
		\ne\varnothing\}.
		\]
		Thus $|I_j|>N/2$ for every
		$j\in J_{\mathrm{exc}}$.
		
		Choose a sufficiently large constant
		$L$, depending only on the dimension
		and the fixed geometric constants.
		
		If
		\[
		|b_i-b_j|\ge LR,
		\]
		then every tube intersecting both
		$P_i$ and $P_j$ has direction within
		angle $O(L^{-1})$ of the unoriented
		line joining $b_i$ and $b_j$.
		
		The angular packing estimate for
		$\delta$-separated directions gives
		\[
		|I_i\cap I_j|
		\le C_n(1+(L\delta)^{-1})^{n-1}.
		\]
		
		Since $N\asymp\delta^{1-n}$,
		we may choose $L$ sufficiently large
		and then $\delta$ sufficiently small
		so that
		\[
		|I_i\cap I_j|\le \frac N6.
		\]
		
		Suppose that three exceptional
		ball centers $b_1,b_2,b_3$
		are pairwise separated by at least $LR$.
		
		By inclusion-exclusion,
		\[
		\begin{aligned}
			\sum_{i=1}^3|I_i|
			&\le
			\left|\bigcup_{i=1}^3 I_i\right|
			+\sum_{1\le i<j\le3}|I_i\cap I_j|\\
			&\le N+3\frac N6\\
			&=\frac{3N}{2}.
		\end{aligned}
		\]
		
		But each $|I_i|>N/2$, so
		\[
		\sum_{i=1}^3|I_i|>\frac{3N}{2},
		\]
		a contradiction.
		
		Consequently, a maximal
		$LR$-separated subfamily of
		exceptional ball centers contains
		at most two centers.
		
		Every exceptional center lies within
		distance $LR$ of one of these centers.
		
		Since the original radius-$R$
		Vitali balls are pairwise disjoint,
		volume packing in the resulting
		two balls of radius $(L+1)R$ gives
		\[
		|J_{\mathrm{exc}}|
		\le 2(L+1)^n.
		\]
		
		For $\delta$ outside the sufficiently
		small range used above, $N$ is bounded.
		Since $q\le N$, the radius $R$ is then
		bounded below. All balls meet
		$F_m\subseteq B(0,1)$, so ordinary
		volume packing supplies a uniform
		bound in this remaining range.
		
		Thus
		\[
		|J_{\mathrm{exc}}|\le C_{n,A}.
		\]
	\end{proof}
		
	Next, we do a elementary estimate. Let $\mathcal P$ be a finite family of nonempty
	localized trimmed pieces. Then 
	\begin{equation} \label{eq:elementary-bound}
		|P| \leq C_nNq^{-1}q^{1/(n-1)} \leq \frac{N}{4},
	\end{equation}
	when $q >> 1.$ We just need to over count the intersections in each tube, and the elementary bound follows. The following is one of our main results. 
	
	\begin{theorem}
		[Existence of an intersection-compatible bijection]
		\label{thm:restricted-bijection}
		
		Let $n\ge3$, $q=2^m\ge8$, $q\le N$,
		and let $T_1^\delta,\ldots,T_N^\delta$
		be the original tubes, with
		$N\asymp\delta^{1-n}$ and pairwise
		$\delta$-separated directions.
		
		Let $B\subseteq\{1,\ldots,N\}$
		be the greedily selected tube indices,
		so that
		\[
		|B|\le\frac{2N}{q}.
		\]
		
		Let $\mathcal V=\{V_j:j\in J\}$
		be a globally disjoint family of
		Vitali balls of radius
		\[
		R=Aq^{-1/(n-1)},\qquad
		\delta\le R\le1.
		\]
		
		Let $\mathcal P$ be a finite family
		of nonempty localized trimmed pieces.
		For every $P\in\mathcal P$, suppose
		there are labels
		$\ell(P)\in B$ and $j(P)\in J$ with
		\[
		P\subseteq
		T_{\ell(P)}^\delta
		\cap G_{\ell(P)}
		\cap F_m\cap V_{j(P)}.
		\]
		
		Assume that
		\[
		P\longmapsto(\ell(P),j(P))
		\]
		is injective.
		
		Define
		\[
		d(P)=
		\#\{r\le N:P\cap T_r^\delta
		\ne\varnothing\}
		\]
		and partition
		\[
		\mathcal P_0=
		\{P\in\mathcal P:d(P)\le N/2\},
		\]
		\[
		\mathcal P_{\mathrm{exc}}=
		\{P\in\mathcal P:d(P)>N/2\}.
		\]
		
		Suppose that
		\[
		|\mathcal P_0|\le\frac N4.
		\]
		
		Then
		\[
		|\mathcal P_{\mathrm{exc}}|
		\le C_{n,A}|B|
		\lesssim\frac Nq.
		\]
		
		Moreover, after adjoining
		$M=|\mathcal P_{\mathrm{exc}}|$
		auxiliary tubes outside $B(0,2)$,
		there exists an injective map
		\[
		\sigma:\mathcal P\longrightarrow
		\{1,\ldots,N+M\}\setminus B
		\]
		such that
		\[
		P\cap T_{\sigma(P)}^\delta
		=\varnothing
		\qquad(P\in\mathcal P).
		\]
		
		Consequently, there exist sets
		\[
		B\subseteq B'\subseteq D
		\]
		and a bijection
		\[
		\Phi:B\sqcup\mathcal P\longrightarrow D
		\]
		fixing every index in $B$, with
		\[
		|B'|=|B|+M
		\lesssim\frac Nq.
		\]
	\end{theorem}
	
	\begin{proof}
		By Lemma~\ref{lem:exceptional-vitali},
		the number of Vitali balls containing
		at least one exceptional piece is at
		most $C_{n,A}$.
		
		For each such ball, the injectivity of
		\[
		P\longmapsto(\ell(P),j(P))
		\]
		implies that at most $|B|$ pieces
		are assigned to that ball.
		
		Therefore,
		\[
		M=|\mathcal P_{\mathrm{exc}}|
		\le C_{n,A}|B|
		\le 2C_{n,A}\frac Nq.
		\]
		
		We first construct representatives
		for the ordinary pieces.
		
		For $P\in\mathcal P_0$, define
		\[
		A(P):=
		\{r\in\{1,\ldots,N\}\setminus B:
		P\cap T_r^\delta=\varnothing\}.
		\]
		
		By definition of $d(P)$,
		\[
		|A(P)|\ge N-d(P)-|B|.
		\]
		
		Since $d(P)\le N/2$ and
		\[
		|B|\le\frac{2N}{q}\le\frac N4,
		\]
		we have
		\[
		|A(P)|\ge\frac N4.
		\]
		
		For every nonempty
		$S\subseteq\mathcal P_0$,
		\[
		\begin{aligned}
			\left|\bigcup_{P\in S}A(P)\right|
			&\ge\max_{P\in S}|A(P)|\\
			&\ge\frac N4\\
			&\ge|S|.
		\end{aligned}
		\]
		
		Thus Hall's marriage condition holds.
		
		There is consequently an injection
		\[
		\sigma_0:\mathcal P_0
		\hookrightarrow\{1,\ldots,N\}\setminus B
		\]
		satisfying
		\[
		P\cap T_{\sigma_0(P)}^\delta
		=\varnothing.
		\]
		
		Next, introduce $M$ additional labels
		\[
		E=\{N+1,\ldots,N+M\}.
		\]
		
		Choose the corresponding auxiliary
		tubes outside $B(0,2)$.
		They are disjoint from every
		$P\subseteq F_m\subseteq B(0,1)$.
		
		Since $|E|=|\mathcal P_{\mathrm{exc}}|$,
		choose a bijection
		\[
		\sigma_{\mathrm{exc}}:
		\mathcal P_{\mathrm{exc}}\longrightarrow E.
		\]
		
		Then
		\[
		P\cap T_{\sigma_{\mathrm{exc}}(P)}^\delta
		=\varnothing
		\]
		for every exceptional piece.
		
		Combining the two maps gives an injection
		\[
		\sigma:\mathcal P\longrightarrow
		\{1,\ldots,N+M\}\setminus B.
		\]
		
		Define
		\[
		B':=B\cup E
		\]
		and
		\[
		D:=B\cup\sigma(\mathcal P).
		\]
		
		Since $\sigma_{\mathrm{exc}}$ is onto $E$,
		we have
		\[
		B\subseteq B'\subseteq D.
		\]
		
		Furthermore,
		\[
		\begin{aligned}
			|B'|
			&=|B|+M\\
			&\le(1+C_{n,A})|B|\\
			&\le2(1+C_{n,A})\frac Nq.
		\end{aligned}
		\]
		
		Finally, define
		\[
		\Phi:B\sqcup\mathcal P\longrightarrow D
		\]
		by
		\[
		\Phi(\ell)=\ell\quad(\ell\in B),
		\qquad
		\Phi(P)=\sigma(P)\quad(P\in\mathcal P).
		\]
		
		Since $\sigma$ is injective and its
		image is disjoint from $B$, this
		map is bijective.
		
		This completes the proof.
	\end{proof}
	
	Let
	\[
	\mathcal A
	:=
	\{(\ell,j):\ell\in B',\ 1\le j\le m_\ell\}
	\]
	be the set of selected tube-ball pairs.
	Redefine the bijective re-indexing for the Vitali- balls in accordance with the theorem \ref{thm:restricted-bijection}:
	\[
	\kappa:\mathcal A\longrightarrow D,
	\qquad
	(\ell,j)\longmapsto \kappa(\ell,j),
	\]
	with inverse
	\[
	\kappa^{-1}(\kappa_0)
	=
	(\ell(k_0),j(\kappa_0)),
	\qquad \kappa_0 \in D'
	\]
	In particular,
	\[
	b_{\kappa(\ell,j)}=b_{\ell,j},
	\qquad
	b_{\kappa_0}=b_{\ell(\kappa_0),j(\kappa_0)}.
	\] 
	So we have
	\[
	\begin{aligned}
	\#\{\ell \in B \sqcup \sigma \mathcal P: \Phi(\ell) = r\} 
	= 1, 
	\end{aligned}
	\]
	for all $r \in  B \sqcup \mathcal P,$
	using the mapping given by the theorem \ref{thm:restricted-bijection},
	We extend the bijection $\Phi$ to the original tube
	indices $\ell=1,\ldots,N$ by the incidence rule
	\[
	\widehat{\Phi}(\ell)=r
	\quad\Longleftrightarrow\quad
	T_\ell^\delta\cap P_r\cap G_r\ne\varnothing,
	\qquad r\in D.
	\]
	Here $T_\ell^\delta$ is an original, untrimmed tube,
	whereas $P_r\subseteq F_m\cap G_r$ is a retained
	localized piece.
	The uniqueness of the target $r$ follows from the greedy selection \ref{thm:greedy-selection}.  
	Thus, by double counting,
	\[
	\begin{aligned}
		q|\mathcal P|
		&\leq 
	    \sum_{r \in \sigma \mathcal P}
		\#\{\ell \in \{1,\ldots,N+M\}: \Phi(\ell) = r\}  \\
		&= 
		\sum_{r = 1}^{N+M}
		\#\{\ell \in \sigma \mathcal P: \Phi(\ell) = r\}  \\
		&\leq N + M \\
		&\lesssim N
		,
	\end{aligned}
	\]
	So we have $|D| \lesssim 2^{-m}N.$
	So we have 
	\begin{equation}
		|F_m| \lesssim \sum_ {j=1}^{|D|}|T^\delta_j(a_j) \cap B(b_{j},3C2^{-m/(n-1)}) \cap G_j \cap F_m| \lesssim \delta^{n-1}|D|2^{-m/(n-1)} \lesssim 2^{-mn/(n-1)} 
	\end{equation}
	So the Theorem~\ref{thm:main} holds after a simple calculation from \eqref{eq:final1}, \eqref{eq:final2},\eqref{eq:final3}, and above:
	\[
	\ln N \sim \sum_m 2^{-m(2-n/(n-1))} 2^{m}2^{-m/(n-1)}
	\gtrsim \left\| \sum_{i=1}^{N} \mathbf{1}_{T_i(a_i)}(x)
	\right\|_{n/(n-1)}^{n/(n-1)}.
	\] 
	
	\section{Appendixes}
	\appendix
	Some of this results are used in the Lean project \cite{Aspegren2026Lean}.
	We establish three elementary results on selected-incidence counting: a weighted selection comparison, an equivalence between lower intersection incidence and representative-map surjectivity for nonempty disjoint pieces, and a surjectivity criterion for representative maps via composition. 
	
	Finite incidence arguments arise naturally when a large family of sets is
	reduced to a smaller family while retaining information about their
	overlaps. A common situation is that geometric objects, such as tubes,
	are indexed by $j\in\{1,\ldots,N\}$, while the intersections or localized
	pieces associated with these objects carry a second index.  The resulting
	double-indexed union is conveniently reindexed by a single finite set,
	but a change of indices must not be confused with a reduction in the
	number of pieces.
	
	This note isolates three elementary principles behind such arguments.
	The first is a weighted \emph{selection comparison}: provided every
	piece remains incident to at least one selected object and original
	incidence multiplicities are bounded, the selected family controls the
	total incidence weight. The second identifies the precise meaning of
	lower intersection incidence for nonempty, pairwise-disjoint pieces:
	it is equivalent to surjectivity of a map selecting a representative
	piece. The third relates this surjectivity to a composition: if the
	piece map is injective and the composite representative-to-piece-set
	map is onto the range of the piece map, then the representative map is
	onto the piece labels.
	
	The arguments developed in this paper are based on
	finite incidence counting, combinatorial inequalities,
	and properties of injective and surjective mappings.
	For general background on combinatorial counting
	methods, we refer to \cite{vanLintWilson}.
	The geometric setting of doubly indexed disjoint
	families is related to the broader framework of
	discrete geometry; see \cite{Matousek}.
	
	\section{Selection comparison}
	
	Let $\calP$ be a finite set of labels for measurable pieces. The
	measure-theoretic structure is not used below; only a nonnegative weight
	$c_P$ associated with each $P\in\calP$ is needed.
	
	\begin{theorem}[Selection comparison]\label{thm:comparison}
		For each $j\in\{1,\ldots,N\}$, let $\calP_j\subseteq\calP$ be an
		original incidence family. Assign weights $c_P\geq 0$ and put
		\[
		w_j:=\sum_{P\in\calP_j}c_P,
		\qquad
		r_P:=\#\bigl\{j\in\{1,\ldots,N\}:P\in\calP_j\bigr\}.
		\]
		Suppose that, for some $q\geq1$ and $C>0$,
		\[
		r_P\leq Cq\qquad(P\in\calP).
		\]
		Let $A\subseteq\{1,\ldots,N\}$ be a selected family such that every
		piece belongs to at least one selected incidence family:
		\[
		s_P:=\#\{j\in A:P\in\calP_j\}\geq1
		\qquad(P\in\calP).
		\]
		Then
		\begin{equation}\label{eq:selection-comparison}
			\boxed{\displaystyle
				\sum_{j=1}^{N}w_j\leq Cq\sum_{j\in A}w_j.}
		\end{equation}
	\end{theorem}
	
	\begin{proof}
		For each $P\in\calP$, the assumptions give
		$r_P\leq Cq\leq Cq\,s_P$. As $c_P\geq0$, finite double counting yields
		\begin{align*}
			\sum_{j=1}^{N}w_j
			&=\sum_{j=1}^{N}\sum_{P\in\calP_j}c_P
			=\sum_{P\in\calP}r_Pc_P\\
			&\leq Cq\sum_{P\in\calP}s_Pc_P
			=Cq\sum_{j\in A}\sum_{P\in\calP_j}c_P
			=Cq\sum_{j\in A}w_j.
		\end{align*}
	\end{proof}
	
	\noindent
	In particular, if each piece belongs to \emph{exactly one} selected
	family, then $s_P=1$; however, the proof only requires $s_P\geq1$.
	An upper bound $s_P\leq1$ alone does not establish
	\eqref{eq:selection-comparison}.
	
	\section{Lower intersection incidence}
	
	Let $P$, $J$ and $X$ be sets, and let
	\[
	E:P\longrightarrow\mathcal P(X),\qquad
	\rep:J\longrightarrow P.
	\]
	Here $E(p)$ is the piece indexed by $p$, while $\rep(j)$ selects a
	piece label for the representative index $j$. Define
	\begin{equation}\label{eq:lower-incidence}
		\LowerIncidence(E,\rep)
		\quad\Longleftrightarrow\quad
		\forall p\in P\ \exists j\in J:\ 
		E(p)\cap E(\rep(j))\ne\varnothing.
	\end{equation}
	
	\begin{theorem}[Lower incidence is equivalent to representative
		surjectivity]\label{thm:lower-surjectivity}
		Suppose that each $E(p)$ is nonempty and that the pieces are pairwise
		disjoint in the indexed sense:
		\[
		E(p)\ne\varnothing\quad(p\in P),\qquad
		p\ne p'\ \Longrightarrow\ E(p)\cap E(p')=\varnothing.
		\]
		Then
		\begin{equation}\label{eq:iff-surj}
			\boxed{
				\LowerIncidence(E,\rep)
				\quad\Longleftrightarrow\quad
				\rep:J\twoheadrightarrow P.}
		\end{equation}
		Furthermore, for every $j\in J$ and $p\in P$,
		\begin{equation}\label{eq:exact-incidence}
			E(p)\cap E(\rep(j))\ne\varnothing
			\quad\Longleftrightarrow\quad p=\rep(j).
		\end{equation}
	\end{theorem}
	
	\begin{proof}
		If $p=\rep(j)$, then
		$E(p)\cap E(\rep(j))=E(p)\ne\varnothing$.
		Conversely, if this intersection is nonempty, pairwise disjointness
		forces $p=\rep(j)$. This establishes \eqref{eq:exact-incidence}.
		
		Suppose $\LowerIncidence(E,\rep)$. For arbitrary $p\in P$,
		choose $j\in J$ with
		$E(p)\cap E(\rep(j))\ne\varnothing$. By
		\eqref{eq:exact-incidence}, $p=\rep(j)$; hence $\rep$ is surjective.
		
		Conversely, if $\rep$ is surjective, then for every $p\in P$ we can
		choose $j\in J$ with $\rep(j)=p$. Since $E(p)$ is nonempty,
		$E(p)\cap E(\rep(j))=E(p)\ne\varnothing$. This proves lower incidence.
	\end{proof}
	
	\noindent
	Equation \eqref{eq:exact-incidence} also gives an exact row count:
	when $P$ is finite, for every fixed $j\in J$,
	\[
	\#\{p\in P:E(p)\cap E(\rep(j))\ne\varnothing\}=1.
	\]
	
	
	\section{Surjectivity of the composed piece map}
	
	The natural codomain for the composition $E\circ\rep$ is not all of
	$\mathcal P(X)$, but the range
	\[
	\range(E):=\{E(p):p\in P\}\subseteq\mathcal P(X).
	\]
	Thus the induced composite is
	\[
	\widetilde{E\circ\rep}:J\longrightarrow\range(E),
	\qquad j\longmapsto E(\rep(j)).
	\]
	
	\begin{theorem}[Surjective composition with injective outer
		map]\label{thm:composite}
		Let $E:P\to\mathcal P(X)$ be injective, and let
		$\rep:J\to P$ be any map. If the induced composition
		\[
		\widetilde{E\circ\rep}:J\longrightarrow\range(E)
		\]
		is surjective, then $\rep$ is surjective.
	\end{theorem}
	
	\begin{proof}
		Fix $p\in P$. Since $E(p)\in\range(E)$, surjectivity of
		$\widetilde{E\circ\rep}$ supplies $j\in J$ such that
		\[
		E(\rep(j))=E(p).
		\]
		Injectivity of $E$ gives $\rep(j)=p$. Thus every $p\in P$ lies in
		the range of $\rep$, which proves its surjectivity.
	\end{proof}
	
	\noindent
	\textbf{Connection between the three theorems.}
	Nonempty, pairwise-disjoint indexed pieces make $E$ injective: if
	$E(p)=E(p')$ and $p\ne p'$, the two nonempty equal sets would be
	disjoint, a contradiction. Hence, in that setting, surjectivity of
	$\widetilde{E\circ\rep}$ implies surjectivity of $\rep$ by
	Theorem~\ref{thm:composite}, and then lower incidence by
	Theorem~\ref{thm:lower-surjectivity}. When these intersections encode the
	selected incidence families of Theorem~\ref{thm:comparison}, the
	resulting column coverage supplies its hypothesis $s_P\geq1$.
	
	
	
	\section*{Acknowledgments}
	
	The author used OpenAI's ChatGPT, and Anthropic's Claude, to assist with
	mathematical discussions, proof review, LaTeX
	preparation, and Lean 4 formalization.
	The mathematical arguments and responsibility
	for their correctness remain with the author.
	
	\begin{thebibliography}{3}
		
		\bibitem{Aspegren2026Lean}
		J. Aspegren,
		\emph{KakeyaProject: Formalization of the Kakeya maximal function argument in Lean 4},
		GitHub repository, 2026.
		\url{https://github.com/joas165/KakeyaProject}.
		
		
		\bibitem{Kroc}
		
		{\scshape Kroc, E.} 
		
		The Kakeya problem. 
		
		Available at 
		
		\url{http://ekroc.weebly.com/uploads/2/1/6/3/21633182/mscessay-final.pdf}.
		
		\bibitem{vanLintWilson}
		J. H. van Lint and R. M. Wilson,
		\textit{A Course in Combinatorics},
		2nd ed., Cambridge University Press, 2001.
		
		\bibitem{Matousek}
		J. Matou\v{s}ek,
		\textit{Lectures on Discrete Geometry},
		Graduate Texts in Mathematics, Vol. 212,
		Springer, New York, 2002.
		
		\bibitem{Tao99}
		
		{\scshape Tao, Terence.}
		
		Lecture notes for harmonic analysis course. 
		
		University of California, Los Angeles, 1999. 
		
		Available at 
		
		\url{http://math.ucla.edu/~tao/254b.1.99s/}.
		
		
		
		
	\end{thebibliography}
	
	
	\end{document}
	
	\backmatter
	
	\section*{Declarations}
	
	\bmhead{Funding}
	
	The author received no funding for this work.
	
	\bmhead{Conflict of interest/Competing interests}
	
	The author declares no conflict of interest.
	
	\bmhead{Ethics approval and consent to participate}
	
	Not applicable.
	
	\bmhead{Consent for publication}
	
	Not applicable.
	
	\bmhead{Data availability}
	
	Not applicable.